\section{Bosquejo del Método}

\subsection{Población}

La presente investigación tomará como enfoque principal la población del
municipio de Culiacán, Sinaloa, México. De acuerdo con el último Censo de
Población y Vivienda realizado por el Instituto Nacional de Estadística y
Geografía \parencite[p.~23]{inegi2021}, Culiacán registra una población total de
1,003,530 habitantes, lo que representa el 33.2\% de la población total del
estado de Sinaloa.

Dentro de este universo, se delimitó como población de estudio a los habitantes
de 27 años de edad. Este grupo resulta de particular interés debido a que ha
transitado por las distintas etapas del comercio moderno: desde la adquisición
de productos exclusivamente a través de catálogos o tiendas físicas, hasta la
adopción de las compras en línea, cada una con estrategias de marketing propias.
Dicha trayectoria les otorga una perspectiva amplia frente al consumo de
cualquier tipo de producto. Asimismo, por tratarse de adultos en plena vida
laboral —y en la mayoría de los casos con los estudios concluidos—, es razonable
suponer que ya han adquirido los tipos de productos en los que se centra esta
investigación: laptops y smartphones. A esto se suma que una persona en esta
etapa de vida cuenta con mayor criterio, autonomía económica y experiencia
acumulada en decisiones de compra.

De acuerdo con el Censo de Población y Vivienda 2020, la población de Culiacán
correspondiente al rango de 27 años o más asciende a 578,109 habitantes. Este
dato es congruente con la edad mediana registrada por el propio INEGI para ese
año, que fue de 29 años; en efecto, el 50\% de la población total del municipio
equivale a:

\[
\frac{1{,}003{,}530}{2} = 501{,}765 \text{ habitantes}
\]

lo que confirma que más de la mitad de la población se concentra en edades
cercanas al rango de interés.

Con el fin de obtener una estimación más precisa para el año 2026, se incorporó
al cálculo la población que, en 2020, se encontraba en el rango de 20 a 24 años,
puesto que dichas personas habrán alcanzado los 27 años o más transcurridos seis
años. Al sumar este grupo al total anterior, se obtiene una base poblacional
ampliada de:

\[
578{,}109 + 88{,}495 = 666{,}604 \text{ habitantes}
\]

Además, para proyectar esta cifra a 2026, se descontó la mortalidad estimada del
periodo. Tomando como referencia un promedio de 18,902 defunciones anuales
registradas en Sinaloa a partir de 2020 \parencite{inegi_mortalidad}, y
considerando que Culiacán concentra el 33.2\% de la población estatal, le
corresponden aproximadamente:

\[
18{,}902 \times 0.332 \approx 6{,}275 \text{ defunciones anuales}
\]

Descontando estas defunciones acumuladas durante el periodo 2020--2026, la
población estimada en el rango de 27 años o más para el presente año es de:

\[
666{,}604 - (6{,}275 \times 6) = 666{,}604 - 37{,}650 \approx 628{,}954 
\text{ habitantes}
\]

\subsection{Muestra}

Tomando los valores anteriormente detallados para la población a estudiar, se
procedió a calcular el tamaño de muestra necesario. Empleando la fórmula
propuesta por \textcite[p.~75]{cochran1977}:
\[
n = \frac{\dfrac{t^2 PQ}{d^2}}{1 + \dfrac{1}{N}\left(\dfrac{t^2 PQ}{d^2} - 1\right)}
\]

\noindent con un nivel de confianza del 95\% ($t = 1.96$), una proporción
esperada de $P = 0.5$, y un margen de error de $d = 0.05$, se obtuvo:

\[
n = \frac{\dfrac{(1.96)^2 (0.5)(0.5)}{(0.05)^2}}{1 + \dfrac{1}{1{,}003{,}530}
\left(\dfrac{(1.96)^2 (0.5)(0.5)}{(0.05)^2} - 1\right)}
= \frac{384.16}{1 + \dfrac{383.16}{1{,}003{,}530}} \approx 384 \textrm{ encuestados}
\]

Posteriormente, aplicando la asignación proporcional por estrato
\parencite[p.~91]{cochran1977}, se determinó la muestra correspondiente al
segmento de habitantes de interés con respecto a la población total del
municipio:

\[
n_h = n \cdot \frac{N_h}{N} = 384 \cdot \frac{628{,}954}{1{,}003{,}530} 
\approx 241 \textrm{ encuestados}
\]

Por lo tanto, la muestra definitiva para esta investigación quedó establecida en
\textbf{241 encuestados}, representando de manera proporcional al segmento
poblacional de interés dentro del municipio de Culiacán, Sinaloa.

\subsection{Tipo de Estudio}

La presente investigación adopta un enfoque metodológico mixto. En primer lugar,
es de carácter documental, dado que su sustento teórico está construido a partir de
estudios/experimentos previos que respaldan los efectos positivos propuestos que
la disponibilidad de la información —presentados de forma adecuada— tiene 
en la toma de desiciones del consumidor.

Asimismo, la investigación integra un componente cuantitativo
y cualitativo mediante la aplicación de una encuesta —cuyo instrumento se
presenta en el Anexo A— dirigida a consumidores de dispositivos electrónicos.
Dicho instrumento tiene como propósito recabar la opinión de usuarios reales
respecto a sus hábitos de consumo, su acceso a información al momento de
adquirir un producto, así como la percepción de utilidad que tendría una herramienta
como la propuesta en este proyecto. La muestra representativa está conformada por
241 encuestados que permitirá obtener resultados con un nivel de confianza 
adecuado para validarla relevancia de la aplicación dentro del contexto estudiado.
